
\documentclass[10pt]{article}

\usepackage{cmap}
\usepackage[T2A]{fontenc}
\usepackage[utf8]{inputenc}
\usepackage[english,russian]{babel}
\usepackage{geometry}
\usepackage{hyperref}
\usepackage{enumitem}
\usepackage{array}
\usepackage{fontawesome}
\usepackage{titlesec}
\usepackage{xcolor}

\geometry{
    top=10mm,
    bottom=10mm,
    left=14mm,
    right=14mm
}

\pagenumbering{gobble}
\setlength{\parindent}{0pt}
\setlength{\parskip}{1pt}
\renewcommand{\arraystretch}{0.95}

\hypersetup{
    colorlinks=true,
    urlcolor=blue,
    linkcolor=blue,
    citecolor=black,
    anchorcolor=black
}

\setlist[itemize]{
    noitemsep,
    topsep=0pt,
    partopsep=0pt,
    parsep=0pt,
    itemsep=0pt,
    leftmargin=1.1em
}

\titleformat{\section}{\large\bfseries\raggedright}{}{0em}{}[\titlerule]
\titlespacing{\section}{0pt}{8pt}{5pt}

\newif\ifen
\newif\ifru
\newcommand{\en}[1]{\ifen#1\fi}
\newcommand{\ru}[1]{\ifru#1\fi}

\rutrue
% \entrue

\newcommand{\dates}[2]{\footnotesize #2\\[-1pt]\footnotesize #1}

\newcommand{\workentry}[3]{
\begin{tabular}{@{}p{0.125\linewidth}|p{0.82\linewidth}@{}}
\raggedleft #1 & \raisebox{-1.5pt}[0pt][0pt]{\textbf{#2}} \\[-12pt]
& #3
\end{tabular}\vspace{1mm}
}

\begin{document}

\begin{center}
    {\Huge \textbf{\ru{Никита Зайцев}\en{Nikita Zaytsev}}}
\end{center}

\begin{center}
\small
\faEnvelope\ \href{mailto:nickita-zaitsiev@ya.ru}{nickita-zaitsiev@ya.ru}
\hspace{1cm}
\faPhone\ \href{tel:+79657447577}{+7 965 744-75-77}
\hspace{1cm}
\faPaperPlane\ \href{https://t.me/nickzay}{@nickzay}
\hspace{1cm}
\faGithub\ \href{https://github.com/NickZay}{NickZay}
\hspace{1cm}
\faLinkedin\ \href{https://www.linkedin.com/in/nickzay}{nickzay}
\hspace{1cm}
\textbf{hh} \href{https://hh.ru/resume/e86c2e50ff0bf0d5320039ed1f686758497349}{resume}
\end{center}

\vspace{-4mm}
\section{\ru{О себе}\en{Summary}}

\ru{Аналитик данных с сильным инженерным бэкграундом и более чем четырьмя годами опыта в продуктовых системах Ozon, Яндекса и Сбера. Уверенно работаю с Python, SQL, витринами и большими данными. Понимаю путь данных от источника до аналитики, умею исследовать продуктовые метрики, находить узкие места и автоматизировать расчёты. Рассматриваю позиции data- и продуктового аналитика.}
\en{Data Analyst with a strong engineering background and more than four years of experience with product systems at Ozon, Yandex, and Sber. I work confidently with Python, SQL, data marts, and large-scale data. I understand the data path from source to analytics, can investigate product metrics, identify bottlenecks, and automate calculations. I am considering Data Analyst and Product Analyst positions.}

\vspace{-2mm}
\section{\ru{Опыт работы}\en{Work Experience}}

\workentry
{\dates{\ru{Июнь}\en{June} 2026}{\ru{Август}\en{August} 2026}}
{\ru{Ozon.ru --- Go-разработчик}\en{Ozon.ru --- Go Developer}}
{
\begin{itemize}
    \item \ru{Разрабатывал backend-сервисы и интеграции e-commerce-платформы, работал с потоками событий, хранилищами, метриками и логами}
          \en{Developed backend services and integrations for an e-commerce platform, working with event streams, storage systems, metrics, and logs}
    \item \ru{Стек}\en{Stack}: Go, Python, Kafka, Elasticsearch, PostgreSQL, Redis, K8s, Kibana, Grafana
\end{itemize}
}

\workentry
{\dates{\ru{Сентябрь}\en{September} 2024}{\ru{Май}\en{May} 2026}}
{\ru{Яндекс Плюс --- Frontend-разработчик}\en{Yandex Plus --- Frontend Developer}}
{
\begin{itemize}
    \item \ru{Реализовал допродажу тарифов и опций на главном экране оплаты: продажи дополнительных опций выросли на 1\%}
          \en{Implemented plan and option upsell on the main payment screen, increasing add-on sales by 1\%}
    \item \ru{Реализовал онбординг после оплаты подписки, что повысило вовлечённость пользователей на 5\%}
          \en{Implemented post-payment onboarding, increasing user engagement by 5\%}
    \item \ru{Добавил историю изменений промокампаний и развивал сценарии оплаты и активации подписки}
          \en{Added promo campaign change history and improved subscription payment and activation flows}
    \item \ru{Стек}\en{Stack}: TypeScript, React, Next.js, XState, Redux, Playwright, Node.js
\end{itemize}
}

\workentry
{\dates{\ru{Февраль}\en{February} 2024}{\ru{Сентябрь}\en{September} 2024}}
{\ru{Arenadata (Picodata) --- Rust-разработчик}\en{Arenadata (Picodata) --- Rust Developer}}
{
\begin{itemize}
    \item \ru{Писал и оптимизировал Ansible-сценарии для воспроизводимого развёртывания распределённых приложений}
          \en{Developed and optimized Ansible workflows for reproducible deployment of distributed applications}
    \item \ru{Перевёл автотесты с Lua на Rust: время выполнения сократилось на 10\%, разработка и тестирование были унифицированы на одном языке}
          \en{Migrated automated tests from Lua to Rust, reducing execution time by 10\% and unifying development and testing on a single language}
    \item \ru{Стек}\en{Stack}: Rust, Tarantool, Kafka, Lua, Ansible
\end{itemize}
}

\workentry
{\dates{\ru{Август}\en{August} 2022}{\ru{Сентябрь}\en{September} 2023}}
{\ru{Сбер --- Аналитик данных}\en{Sberbank --- Data Analyst}}
{
\begin{itemize}
    \item \ru{Разрабатывал витрины и расчёты для продукта групп связанных заёмщиков: агрегировал кредитную экспозицию, анализировал использование лимитов и состав групп}
          \en{Developed data marts and calculations for a related-borrower group product, aggregated credit exposure, and analyzed limit utilization and group composition}
    \item \ru{Контролировал качество связей между заёмщиками, выявлял дубли и расхождения между источниками, оптимизировал SQL-расчёты на больших данных}
          \en{Controlled borrower relationship data quality, identified duplicates and discrepancies across sources, and optimized SQL calculations on large datasets}
    \item \ru{Стек}\en{Stack}: SQL, Python, Hadoop, Hive, Spark, Greenplum, ClickHouse
\end{itemize}
}

\workentry
{\dates{\ru{Июль}\en{July} 2021}{\ru{Август}\en{August} 2022}}
{\ru{Яндекс Go --- C++ Backend-разработчик}\en{Yandex Go --- C++ Backend Developer}}
{
\begin{itemize}
    \item \ru{Мигрировал базы данных под новые бизнес-сущности и добавил прогресс достижений водителей: переходы на вкладку достижений выросли на 45\%}
          \en{Migrated databases for new business entities and added driver achievement progress tracking, increasing visits to the achievements tab by 45\%}
    \item \ru{Реализовал кеширование и пакетные API-методы, увеличив максимальную пропускную способность в 10 раз без существенной потери актуальности данных}
          \en{Implemented caching and bulk API endpoints, increasing maximum throughput by 10x without significantly reducing data freshness}
    \item \ru{Стек}\en{Stack}: C++, Python, PostgreSQL, MongoDB, Elasticsearch, Kibana
\end{itemize}
}

\vspace{-4mm}
\section{\ru{Навыки}\en{Skills}}

\begin{tabular}{@{}p{0.20\linewidth}|p{0.75\linewidth}@{}}
\ru{Анализ данных}\en{Data analysis} & SQL, Python, Pandas, NumPy, Statsmodels \\
\ru{Визуализация}\en{Visualization} & DataLens, Matplotlib, Grafana, Kibana \\
\ru{Данные и потоки}\en{Data platforms} & ClickHouse, PostgreSQL, Greenplum, Hadoop, Hive, Spark, Kafka \\
\ru{Аналитические задачи}\en{Analytics} & \ru{витрины, качество данных, продуктовые метрики, оптимизация SQL}\en{data marts, data quality, product metrics, SQL optimization} \\
\ru{Разработка}\en{Development} & Go, TypeScript, React, Java, Rust, C++ \\
\ru{Инструменты}\en{Tools} & Git, Docker, K8s, Ansible, Playwright \\
\end{tabular}

\vspace{-2mm}
\section{\ru{Образование}\en{Education}}

\begin{tabular}{@{}p{0.125\linewidth}|p{0.82\linewidth}@{}}
\raggedleft 2019 -- 2023 &
\raisebox{-1.5pt}[0pt][0pt]{\textbf{\ru{МФТИ}\en{MIPT}} --- \ru{Бакалавр, Прикладная математика и информатика}
\en{Bachelor's degree in Applied Mathematics and Computer Science}} \\
\raggedleft 2023 -- 2025 &
\raisebox{-1.5pt}[0pt][0pt]{\textbf{\ru{МФТИ}\en{MIPT}} --- \ru{Магистратура, Прикладная математика и информатика}
\en{Master's degree in Applied Mathematics and Computer Science}} \\
\raggedleft 2023 -- 2025 &
\raisebox{-1.5pt}[0pt][0pt]{\textbf{\ru{ШАД}\en{YSDA}} --- \ru{Анализ данных}
\en{Data Analysis}}
\end{tabular}

\vfill
\begin{center}
\large
\href{https://github.com/NickZay/CV}{\ru{Актуальная версия этого резюме}\en{Current version of this CV}}
\end{center}

\end{document}
